\documentclass{article}
\usepackage{amsmath,amssymb,mathtools,hyperref}
\usepackage[margin=1in]{geometry}
\allowdisplaybreaks
\begin{document}

\section*{GP learning-based platooning MPC. Approved finite-horizon edition of (17)}
\begin{equation}
\begin{aligned}
\min_{\mathbb V}\;&
\sum_{\mathbf{n_a}=1}^{N_a}\sum_{i=k}^{k+N-1}
 \left\|\mathrm{acc}_{i|k}^{\mathbf{n_a}}\right\|_R^2
+\sum_{i=k}^{k+N-1}
 \left\|v_{i+1|k}^{1}-v_{i+1|k}^{\mathrm{ref}}\right\|_{Q_1}^2\\
&+\sum_{\mathbf{n_a}=2}^{N_a}\sum_{i=k}^{k+N-1}
 \left\|v_{i+1|k}^{\mathbf{n_a}}-v_{i+1|k}^{\mathbf{n_a}-1}\right\|_{Q_2}^2 .
\end{aligned}\tag{17a}
\end{equation}
With the printed decision collection
\[
\mathbb V=\left\{v_{i|k}^{1},v_{i|k}^{\mathbf{n_a}},v_{i|k}^{H},
 p_{i|k}^{\mathbf{n_a}},\mu_{i|k}^{p^H},\Sigma_{i|k}^{p^H},
 \mathrm{acc}_{i|k}^{\mathbf{n_a}}\right\},
\]
subject to
\begin{align}
v_{i+1|k}^{\mathbf{n_a}}&=v_{i|k}^{\mathbf{n_a}}+T\,\mathrm{acc}_{i|k}^{\mathbf{n_a}},
& p_{i+1|k}^{\mathbf{n_a}}&=p_{i|k}^{\mathbf{n_a}}+T\,v_{i|k}^{\mathbf{n_a}},\tag{17b}\\
v_{i|k}^{H}&=f\left(v_{i-1:i-4|k}^{H},v_{i-1:i-4|k}^{N_a}\right),\tag{17c}\\
\mu_{i+1|k}^{p^H}&=\mu_{i|k}^{p^H}+T\,v_{i|k}^{H}
 +T\,\mu^d\left(v_{i-1|k}^{H},v_{i-1|k}^{N_a}\right),\tag{17d}\\
\Sigma_{i+1|k}^{p^H}&=\Sigma_{i|k}^{p^H}
 +T^2\Sigma^d\left(v_{i-1|k}^{H},v_{i-1|k}^{N_a}\right),\tag{17e}\\
p_{i|k}^{\mathbf{n_a}-1}-p_{i|k}^{\mathbf{n_a}}&\geq\Delta,\tag{17f}\\
p_{i|k}^{N_a}-\mu_{i|k}^{p^H}&\geq\Delta+\phi^{-1}(p_{\mathrm{def}})
 \sqrt{\Sigma_{i|k}^{p^H}},\tag{17g}\\
v_{\min}\leq v_{i|k}^{\mathbf{n_a}}&\leq v_{\max},
&\mathrm{acc}_{\min}\leq\mathrm{acc}_{i|k}^{\mathbf{n_a}}&\leq\mathrm{acc}_{\max}.\tag{17h}
\end{align}

\paragraph{Scope, roles, and printed indexing (pp.\ 19, 22--23).}
The longitudinal platoon consists of $N_a$ AVs, denoted $A^{\mathbf{n_a}}$, followed by one HV, denoted $H$. Section 3.1 prints the vehicle-index description as
\[
\mathbf{n_a}=\{1,2,\ldots,N_a\}.
\]
Here $N_a$ is the number of AVs, while $\mathbf{n_a}$ is used as a component index in the equations. The leader is AV $1$, and AV $N_a$ immediately precedes the HV. Positions, velocities, and accelerations are denoted by $p$, $v$, and $\mathrm{acc}$, respectively. The quantities $\mu^{p^H}$ and $\Sigma^{p^H}$ are the HV position mean and variance. The current discrete time is $k$; $N$ is the prediction horizon. The reference velocity $v^{\mathrm{ref}}$, velocity and acceleration limits, safe distance $\Delta$, required satisfaction probability $p_{\mathrm{def}}$, model parameters, and trained GP are fixed data, not additional elements of $\mathbb V$. The cost weights and sampling-time premise are
\[
Q_1>0,\qquad Q_2>0,\qquad R>0,\qquad 0<T\ll1.
\]
The paper leaves the weighted norms in (17a) abstract. In the synthetic family all velocities and accelerations are scalar and $\|z\|^2_Q=Qz^2$, with $Q_1,Q_2,R,T>0$.

\paragraph{Prediction interval and source discrepancy.}
This explicitly approved execution edition uses $N$ controls at $i=k,\dots,k+N-1$ and $N+1$ state snapshots at $i=k,\dots,k+N$. The two velocity terms of (17a) evaluate the $N$ future snapshots at $i+1=k+1,\dots,k+N$. The published (17a) instead ends both velocity sums at $i=k+N$, introducing an additional speed at $k+N+1$ after the last control. The authors' implementation allocates $N$ controls and $N+1$ states. Selecting that horizon changes (17a) as displayed here. It remains possible to supply a separate out-of-horizon reference sample with the historical input. Its value has no effect on this objective.

Apply (17b)--(17e) on $k,\dots,k+N-1$. Apply the speed limits and (17f)--(17g) on $k,\dots,k+N$ and the acceleration limits on $k,\dots,k+N-1$. Indices in (17b), speed limits, and control limits range over $1,\dots,N_a$; (17f) ranges over $2,\dots,N_a$. All required current AV positions and speeds are observations. The initial human position mean is the observed human position and its variance is zero. The human speed and the last AV's speed are observed at the four times $k-4,\dots,k-1$. The observed current human speed agrees with (17c) on this instance family. The GP queried at transition $i$ receives the pair of speeds at $i-1$ and contributes its posterior variance to (17e).

The author manuscript's earlier equation and comments document initialization from observed position and four velocity lags. The authors' Matlab code documents the dimensions of the finite horizon. This edition retains the paper's $R,Q_1,Q_2$ and posterior variance. The author software also evaluates smoothness terms and propagates a standard deviation output; those distinct formulas are recorded separately in the source audit. The reported optimization concerns one receding-horizon update.

\section*{Nominal HV model defining \texorpdfstring{$f$}{f} (pp.\ 8--9)}
The source constructs the nominal model from
\begin{equation}
G_H(s)\approx K\frac{1+T_zs}{1+2\gamma T_ws+T_w^2s^2}e^{-T_ds}
 =\frac{\dot P_H(s)}{\dot P_N(s)}.\tag{1}
\end{equation}
Here $s$ is the Laplace-domain variable, $K$ is the gain, $T_z$ the zero time constant, $\gamma$ the damping coefficient, and $T_d$ the human response delay. The source calls $T_w$ the system's natural frequency. The transforms $\dot P_H(s)$ and $\dot P_N(s)$ correspond to $v_k^H$ and $v_k^{N_a}$. A second-order Pad\'e approximation of the delay and discretization yield
\begin{equation}
\begin{aligned}
v_k^H={}&-c_1v_{k-1}^H-c_2v_{k-2}^H-c_3v_{k-3}^H-c_4v_{k-4}^H\\
&+b_1v_{k-1}^{N_a}+b_2v_{k-2}^{N_a}+b_3v_{k-3}^{N_a}+b_4v_{k-4}^{N_a}\\
={}&f\left(v_{k-1:k-4}^H,v_{k-1:k-4}^{N_a}\right).
\end{aligned}\tag{2}
\end{equation}
The $c_i$ are constant autoregressive coefficients and the $b_i$ constant exogenous-input coefficients, with lag index $i=1,2,3,4$. They are determined by discretizing $G_H(s)$; no numerical values or explicit coefficient-conversion formulas are printed here. The cited discretization methodology in Wang et al.\ (2024) is unverified in the supplied material. The linked ARX+GP construction is
\begin{align}
v_k^H&=\sum_{i=1}^{4}-c_i v_{k-i}^H
       +\sum_{i=1}^{4}b_i v_{k-i}^{N_a}=f(\cdot),\tag{3a}\\
\widetilde v_k^H&=\overbrace{v_k^H}^{\text{ARX prediction}}
 +\overbrace{g(v_{k-1}^H,v_{k-1}^{N_a})}^{\text{GP-based correction}}.\tag{3b}
\end{align}
Thus $v^H$ is the nominal ARX velocity, $\widetilde v^H$ is the GP-compensated velocity, and $g$ learns discrepancies between measured behavior and the nominal model. Equation (17c) uses the nominal $f$, not the compensated velocity as a replacement for $f$.

\section*{GP definitions linked to (17d)--(17e) (pp.\ 10--12)}
The discrepancy observations and inputs are
\begin{equation}
\widehat v_j^H-v_j^H=g(a_{j-1}),\tag{4}
\end{equation}
\[
a_{j-1}=(v_{j-1}^H,v_{j-1}^{N_a}).
\]
Here $\widehat v_j^H$ is the collected velocity and $v_j^H$ the ARX prediction. The printed dataset is
\[
\mathcal D=\left\{
 a=[a_1,\ldots,a_n]^\top\in\mathbb R^{n_a\times n},\quad
 g=[g_1,\ldots,g_n]^\top\in\mathbb R^{1\times n}
\right\}.
\]
In this GP passage, $n$ is the number of data points and plain $n_a$ is the input dimension, not the bold vehicle index $\mathbf{n_a}$ or the number of AVs $N_a$. Each input $a_i$ has dimension $n_a$-by-$1$. The dataset layout and declared dimensions above are retained as supplied; the transpose notation is not silently changed to impose a different matrix layout.

The zero-mean GP uses the squared-exponential kernel
\begin{equation}
k(a_i,a_j)=\sigma_f^2\exp\left[-\frac12(a_i-a_j)^\top L^{-1}(a_i-a_j)\right].\tag{5}
\end{equation}
The hyperparameters $\sigma_f^2$ and $L\in\mathbb R^{n_a\times n_a}$ are optimized by maximizing the log marginal likelihood of the observed discrepancy data. The paper does not supply that optimization as a further controller subproblem. For a new input $a^*\in\mathbb R^{n_a\times1}$,
\begin{align}
g(a^*)&\sim\mathcal N\bigl(\mu^d(a^*),\Sigma^d(a^*)\bigr),\tag{6}\\
\mu^d(a^*)&=K_{a^*a}\bigl(K_{aa}+\sigma^2I\bigr)^{-1}g,\tag{7a}\\
\Sigma^d(a^*)&=K_{a^*a^*}
 -K_{a^*a}\bigl(K_{aa}+\sigma^2I\bigr)^{-1}K_{aa^*}.\tag{7b}
\end{align}
The accompanying Gram-matrix definitions are
\[
[K_{aa}]_{ij}=k(a_i,a_j),\qquad
[K_{aa^*}]=k(a_j,a^*),\qquad
K_{a^*a}=(K_{aa^*})^\top,\qquad
K_{a^*a^*}=k(a^*,a^*).
\]
The symbol $g$ in (7a) is the dataset's discrepancy-output collection, whereas $g(\cdot)$ denotes the GP function. The symbol $\sigma^2$ in (7) is distinct from the kernel hyperparameter $\sigma_f^2$; its definition is not given in the supplied GP passage and is not added here. No unprinted sign, definiteness, or invertibility restrictions on the GP hyperparameters are imposed.

\section*{System assumptions and moment propagation (pp.\ 19--20)}
The underlying AV dynamics, reproduced by (17b), are
\begin{align}
v_{k+1}^{\mathbf{n_a}}&=v_k^{\mathbf{n_a}}+T\,\mathrm{acc}_k^{\mathbf{n_a}},\tag{9a}\\
p_{k+1}^{\mathbf{n_a}}&=p_k^{\mathbf{n_a}}+T\,v_k^{\mathbf{n_a}}.\tag{9b}
\end{align}
AVs are assumed deterministic and to measure and communicate their states flawlessly:
\[
\Sigma(v_k^{\mathbf{n_a}})=0.
\]
The linked HV state equations are
\begin{align}
\widetilde v_k^H&=v_k^H+g(v_{k-1}^H,v_{k-1}^{N_a}),\tag{10a}\\
p_{k+1}^H&=p_k^H+T\widetilde v_k^H.\tag{10b}
\end{align}
At the current time the source assumes
\[
\Sigma(v_k^H)=0,\qquad \Sigma(p_k^H)=0.
\]
The position moments used in (17d)--(17e) are defined by
\begin{equation}
\mu_{k+1}^{p^H}=\mu_k^{p^H}+Tv_k^H
 +T\mu^d(v_{k-1}^H,v_{k-1}^{N_a}),\tag{11}
\end{equation}
with the printed initialization
\[
\mu_0^{p^H}=p_k^H,
\]
and
\begin{align}
\Sigma(p_{k+1}^H)
 &=\Sigma\bigl(p_k^H+Tv_k^H+Tg(v_{k-1}^H,v_{k-1}^{N_a})\bigr),\tag{12a}\\
 &=\Sigma_k^{p^H}+T^2\Sigma_{k-1}^d,\tag{12b}
\end{align}
with
\[
\Sigma_0^{p^H}=0.
\]
The paper expressly states that (12b) overlooks the covariance between $p_k^H$ and $v_k^H$. The zero-subscript initializations and their current-time right-hand side are retained; no new horizon-indexed initialization is substituted. Equation (17e) gives the explicit GP-input arguments of the variance increment appearing in (12b).

\section*{Safety-constraint derivation and selected replacement (pp.\ 20--22)}
The following rows record the source derivation of (17g); they are not additional constraints imposed alongside (17f)--(17g). The preceding AV-spacing prose uses the strict inequality
\[
p_k^{\mathbf{n_a}-1}-p_k^{\mathbf{n_a}}>\Delta.
\]
Equation (17f), retained as the operative row, instead prints $\geq$. The strict prose condition is not also enforced.

For the last AV and HV the source starts with
\begin{equation}
\Pr\bigl(p_k^{N_a}-(p_k^H+\Delta)>\Delta_{\mathrm{ext}}\bigr)
 \geq p_{\mathrm{def}},\tag{13}
\end{equation}
where $\Delta$ is the predefined safe distance, $p_{\mathrm{def}}$ is the required satisfaction probability, and
\[
\Delta_{\mathrm{ext}}\geq0.
\]
The half-space construction is
\[
\mathcal X^{\mathrm{hs}}:=\{x\mid h^\top x\leq b\},\qquad
h\in\mathbb R^n,\quad b\in\mathbb R,
\]
\begin{equation}
\mathcal X^{\mathrm{hs}}(\Sigma_i^x):=
\left\{x\;\middle|\;h^\top x\leq b-\phi^{-1}(p_{\mathrm{def}})
 \sqrt{h^\top\Sigma_i^x h}\right\}.\tag{14}
\end{equation}
For this application,
\[
h^\top=\begin{bmatrix}-1&1\end{bmatrix},\qquad
x:=\begin{bmatrix}p_k^{N_a}&p_k^H+\Delta\end{bmatrix}^{\!\top},\qquad
b=-\Delta_{\mathrm{ext}}.
\]
The source describes $\phi^{-1}$ as the inverse cumulative distribution function. Its covariance display, as supplied in the page-21 text, is
\begin{equation}
\Sigma_k^x:=
\begin{bmatrix}\Sigma_k^{p^{N_a}}\\ \Sigma_k^{p^H}+\Delta\end{bmatrix}
=\begin{bmatrix}0&0\\0&\Sigma_k^{p^H}\end{bmatrix}.\tag{15}
\end{equation}
The page-21 image is not supplied, so the left-hand layout of (15) is image-unverified; no correction to that extracted layout is inferred. The explicit right-hand covariance matrix is accompanied by the assumptions of no covariance between the HV position and its leading AV position and
\[
\Sigma_k^{p^{N_a}}=0.
\]
The source's deterministic replacement is
\begin{equation}
p_k^{N_a}-p_k^H\geq\Delta+\Delta_{\mathrm{ext}}
 +\phi^{-1}(p_{\mathrm{def}})\sqrt{\Sigma_k^{p^H}}.\tag{16}
\end{equation}
The source permits $\Delta_{\mathrm{ext}}$ to be set to zero, describing this option as appropriate when a high level of assurance in constraint satisfaction is desired. Equation (17g) selects this zero-extra-distance form,
\[
\Delta_{\mathrm{ext}}=0,
\]
and uses the predicted HV position mean. The strict event in (13) and non-strict replacement in (16) are retained as printed, not repaired or simultaneously enforced. The source's statement that the resulting requirement exceeds $\Delta$ under all conditions is stronger than (17g) at zero variance, where its right-hand side is exactly $\Delta$; no additional strict margin is introduced.

\section*{Computational dependency boundary (pp.\ 11--12, 17--18, 23--25)}
Equations (6)--(7) are the printed posterior definitions referenced by Section 3.3. The paper also specifies a fully independent conditional (FIC) sparse approximation. It describes compressing the training information into inducing variables whose locations are determined by maximizing the GP marginal likelihood by gradient ascent. Section 2.5 states that the sparse model's hyperparameters come from the full GP models and reports 20 inducing points for its implementation. It does not print a sparse-posterior replacement for (7a)--(7b).

Section 3.4 states that the sparse GP is evaluated on the trajectory obtained at the previous MPC time step and that those results are used in (17d)--(17e) to propagate the current mean and variance. It describes a single sparse GP prediction for the entire prediction horizon rather than separate predictions for each horizon time step. No indexed replacement rows specifying this evaluation or trajectory alignment are printed. Consequently, (17d)--(17e) are retained without inventing frozen-input formulas, inducing-point equations, or an alternative optimization problem. These missing specifications concern the sparse approximation. The full-GP edition above evaluates the posterior using the formulas in Sections 2.1 and 3.3. The authors' Matlab repository was separately inspected for horizon dimensions, trained models and implementation differences.

The numerical sampling time, horizon, weights, bounds, initial configuration, and probability on page 25 are simulation-instance settings, not universal premises of (17), and are not imposed on this symbolic model. The nominal-MPC comparison omits GP effects and is a different controller, not an operative replacement in this reconstruction. The selected Section 3.3 model has $N_a$ AVs and one HV; the page-24 statement that the simulated AV platoon contained ``two HVs'' conflicts with that description and the page-25 description of two AVs followed by one HV, and is not used to change the selected model.

\end{document}
